\documentclass[12pt]{article}
\usepackage[T1]{fontenc}
\usepackage[utf8]{inputenc}
\usepackage{lmodern}
\usepackage{textcomp}
\usepackage{enumitem}
\usepackage{geometry}
\geometry{margin=1in}
\begin{document}
\begin{center}
Answer Key - Philosophy - Undergraduate Level Assessment 16 \
\small (Answers are ordered exactly as in the question paper.)
\end{center}

\section*{Multiple Choice}
\begin{enumerate}[label=\arabic*.]
\item \textbf{Answer:} A
\\ \textit{Options:} A) that is irrelevant as it is never wrong to painlessly kill an animal.; B) that is only wrong if the animal is a pet.; C) that is entirely dependent on the species of the animal.; D) that is dependent on the cultural norms of a society.; E) that is ethically unimportant.
\\ \textit{Note:} In Singer's utilitarian perspective, the ethical consideration revolves around the capacity for suffering rather than the act of killing itself, leading to the conclusion that if an animal does not suffer, it is not inherently wrong to kill it.
\item \textbf{Answer:} A
\\ \textit{Options:} A) just war theory's principle of proportionality.; B) none of the above.; C) all of the above; D) moral relativism.; E) moral nihilism.
\\ \textit{Note:} The principle of military necessity in just war theory is closely related to the principle of proportionality, as both seek to ensure that the use of force is limited to what is necessary to achieve legitimate military objectives while avoiding excessive harm to civilians and non-combatants.
\item \textbf{Answer:} E
\\ \textit{Options:} A) is a form of cognitive distortion.; B) will eventually lead to a decline in student enrollment.; C) will encourage students to disregard positive aspects of a person's history.; D) will lead to a decline in university reputation.; E) will only serve to promote intellectual homogeneity on college and universities.
\\ \textit{Note:} Lukianoff and Haidt contend that disinviting speakers who hold controversial views stifles diverse perspectives, ultimately leading to an environment where only a narrow range of ideas is accepted, thereby fostering intellectual homogeneity in academic institutions.
\item \textbf{Answer:} A
\\ \textit{Options:} A) no valid conclusion can be drawn; B) the conclusion must affirm the antecedent; C) the conclusion must be a tautology; D) the conclusion must affirm the consequent; E) the conclusion must be an inverse statement
\\ \textit{Note:} In a conditional syllogism, denying the antecedent of the conditional statement results in a logically invalid argument because it does not provide sufficient evidence to affirm or deny the consequent, leading to no valid conclusion.
\item \textbf{Answer:} A
\\ \textit{Options:} A) Lhf; B) hLf; C) hLh; D) Lfh; E) LfLh
\\ \textit{Note:} The correct answer is Lhf because it accurately represents the relationship of love between Holly and Frances in predicate logic, where "L" denotes the love relation, "h" represents Holly, and "f" represents Frances.
\end{enumerate}

\section*{True / False}
\begin{enumerate}[label=\arabic*.]
\setcounter{enumi}{5}
\item \textbf{Answer:} True
\\ \textit{Options:} A) True; B) False
\\ \textit{Note:} In Jaina traditions, samyak jnana, or right knowledge, is essential for attaining liberation (moksha) from the cycle of rebirth (samsara) because it involves a profound understanding of the nature of reality, the self, and the principles of karma, which together enable an individual to overcome ignorance and achieve spiritual liberation.
\item \textbf{Answer:} True
\\ \textit{Options:} A) True; B) False
\\ \textit{Note:} Carruthers argues that moral standing is not an absolute concept but varies depending on cultural and philosophical contexts, which aligns with the principles of moral relativism.
\item \textbf{Answer:} False
\\ \textit{Options:} A) True; B) False
\\ \textit{Note:} The fallacy of division is false because it incorrectly assumes that what is true for a whole must also be true for its individual parts, rather than recognizing that the characteristics of parts do not necessarily apply to the entirety.
\item \textbf{Answer:} True
\\ \textit{Options:} A) True; B) False
\\ \textit{Note:} Epicurus believes that temperance, or moderation, is essential for achieving a pleasurable life and that it serves as the foundation for all other virtues, as it enables individuals to manage their desires and attain true happiness.
\item \textbf{Answer:} False
\\ \textit{Options:} A) True; B) False
\\ \textit{Note:} Kant's Humanity formulation asserts that individuals must always be treated as ends in themselves and never merely as means to an end, regardless of the potential benefit to the majority.
\end{enumerate}

\section*{Long Form Response}
\begin{enumerate}[label=\arabic*.]
\setcounter{enumi}{10}
\item \textbf{Answer:} Guru Nanak began preaching the message of the divine Name at the age of 30, a period often associated with maturity and self-awareness. This age likely allowed him to draw upon personal experiences and insights gained during his formative years, thereby enriching his philosophical teachings with depth and relevance. His emphasis on the divine Name as a central tenet in Sikhism reflects a mature understanding of spirituality, advocating for a direct relationship with the divine that transcends ritualistic practices. Moreover, this age coincided with a time of significant social and religious upheaval in South Asia, which may have influenced his call for social justice, equality, and communal harmony. Thus, the age at which he began his ministry was pivotal in shaping his approach to spirituality and social philosophy, making his teachings profoundly impactful.
\\ \textit{Note:} correct because it emphasizes that Guru Nanak's age of 30, a time of maturity, allowed him to integrate personal experiences into his teachings, thereby enhancing the depth and relevance of his message about the divine Name and its significance in establishing a direct relationship with the divine.
\item \textbf{Answer:} Peter Singer's argument in "All Animals Are Equal" is fundamentally rooted in the principal principle, which asserts that equal consideration must be given to the interests of all beings capable of suffering. This principle challenges speciesism, the discrimination against beings solely based on their species, by positing that the capacity to suffer, not species membership, should determine moral consideration. Consequently, Singer argues that just as we reject racism and sexism for their discriminatory nature, we must also extend moral equality to animals, advocating for their rights and welfare. The implications of this argument call for a reevaluation of our ethical frameworks, urging society to recognize the inherent value and rights of non-human animals, ultimately leading to more humane treatment and policies. Thus, Singer's application of the principal principle not only reshapes individual attitudes towards animals but also fosters a broader discourse on ethical responsibilities in our interactions with all sentient beings.
\\ \textit{Note:} correct because it accurately identifies the principal principle as the foundation of Singer's argument, emphasizing the need for equal consideration of interests based on the capacity to suffer, thereby challenging speciesism and drawing parallels with other forms of discrimination like racism and sexism.
\end{enumerate}

\vfill
\noindent\textit{End of Answer Key}
\end{document}